\section{Eigene Methoden}
\label{sec:methode}

\subsection{Datensatz}
\label{subsec:datensatz}

Grundlage der Experimente ist derselbe PatternNet-Datensatz wie in der ersten Arbeit~\cite{steidlehenssler2026milestone1}.
Für die hier untersuchten Distillation-Experimente wurde die Stichprobe wie zuvor auf 100 Bilder pro Klasse begrenzt, damit die Ergebnisse direkt vergleichbar bleiben.

\subsection{Vorverarbeitung und Datenaugmentation}
\label{subsec:augmentation}

Die genaue Vorverarbeitung, Datenaugmentation, Normalisierung und Aufteilung des Datensatzes entsprechen der ersten Arbeit~\cite{steidlehenssler2026milestone1} und werden hier nicht erneut im Detail wiederholt.
Für die vorliegenden Distillation-Experimente wurden diese Schritte unverändert übernommen, um die Vergleichbarkeit mit der Vorarbeit zu sichern.

\subsection{Netzarchitektur}
\label{subsec:architektur}

Die Knowledge-Distillation-Experimente bestehen aus zwei Netzen: einem Teacher und einem Student.
Als Teacher wurde das beste Modell aus den Transfer-Learning-Experimenten verwendet, ein ResNet50 mit dem besten dort erzielten Ergebnis von etwa 97\,\%.
Als Student dient ein abgeschwächtes Modell auf Basis des Base\_Net aus der ersten Arbeit~\cite{steidlehenssler2026milestone1}, das zusätzlich um Batch-Normalisierung erweitert wurde.
Damit bleibt der Student deutlich kleiner als der Teacher, erhält aber eine stabilere Trainingsdynamik als das ursprüngliche Basismodell.

Beide Netze geben vor der Softmax-Aktivierung Logits aus.
Diese Logits sind die Grundlage für die Distillation, da nicht nur die harten Klassenlabels, sondern auch die weichen Vorhersagen des Teachers in die Optimierung einfließen.
Das grundlegende Prinzip der Knowledge Distillation folgt dabei der Arbeit von Hinton et al.~\cite{hinton2015distilling}.

\subsection{Knowledge Distillation}
\label{subsec:distillation}

Seien $z_s$ die Logits des Student-Modells und $z_t$ die Logits des Teacher-Modells.
Aus diesen Werten werden temperaturgeskalierte Verteilungen berechnet:
\begin{equation}
  p_s^{(T)} = \operatorname{softmax}\left(\frac{z_s}{T}\right),
  \qquad
  p_t^{(T)} = \operatorname{softmax}\left(\frac{z_t}{T}\right).
\end{equation}
Die Temperatur $T$ steuert die Weichheit der Verteilungen. Höhere Temperaturen glätten die Wahrscheinlichkeiten und machen Nebenklassen sichtbarer, während kleinere Temperaturen die Verteilungen schärfer machen.

Die verwendete Distillation-Loss kombiniert den Kreuzentropie-Term mit einer KL-Divergenz zwischen Teacher- und Student-Verteilungen:
\begin{equation}
  \mathcal{L}_{\mathrm{KD}} = \alpha \, \mathcal{L}_{\mathrm{CE}}(y, z_s)
  + (1-\alpha) \, T^2 \, \mathcal{L}_{\mathrm{KL}}\left(p_t^{(T)}, p_s^{(T)}\right).
\end{equation}
Dabei ist $y$ das Ground-Truth-Label, $\mathcal{L}_{\mathrm{CE}}$ der Kreuzentropie-Term und $\mathcal{L}_{\mathrm{KL}}$ die Kullback-Leibler-Divergenz zwischen den temperaturgeskalierten Verteilungen.
Der Faktor $T^2$ kompensiert dabei den Skalierungseffekt der Temperatur auf die Gradienten und hält den Beitrag des Distillation-Terms in einem vergleichbaren numerischen Bereich.
Der Parameter $\alpha$ gewichtet das Verhältnis zwischen den harten Labels und dem Wissen des Teachers.

\subsection{Trainingsverfahren}
\label{subsec:training}

Alle Läufe wurden unter identischen Bedingungen trainiert, um die verschiedenen Distillationseinstellungen direkt vergleichen zu können.
Als Optimizer wurde Adam mit den Standardparametern der PyTorch-Implementierung verwendet ($\eta = 10^{-3}$; $\beta_1 = 0{,}9$; $\beta_2 = 0{,}999$).
Ein Learning-Rate-Scheduler wurde nicht eingesetzt, sodass die Lernrate über alle Epochen konstant blieb.

Das Training erfolgte über 15~Epochen mit einer Batchgröße von 64.
Für die Distillation wurden die Temperaturen $T \in \{2; 4; 6\}$ und die $\alpha$-Werte $\alpha \in \{{,}1; {,}2; {,}3\}$ untersucht.
Damit ergaben sich insgesamt neun Durchläufe, in denen jede Temperatur mit jedem Alpha-Wert kombiniert wurde.

Die Bewertung erfolgte primär über die Test-Accuracy.
Zusätzlich wurden nach jeder Epoche die Validierungs-Accuracy und der Validierungs-Loss berechnet und mit TensorBoard protokolliert.
Auch die Konfusionsmatrix wurde auf der Validierungsmenge pro Epoche erstellt, um klassenspezifische Fehlklassifikationen sichtbar zu machen.

\subsection{Implementierungsdetails}
\label{subsec:implementierung}

Die technischen Details zur Dataset-Klasse, zur Stratifizierung, zur DataLoader-Konfiguration sowie zum Logging entsprechen ebenfalls weitgehend der ersten Arbeit~\cite{steidlehenssler2026milestone1}.
Für die hier untersuchten Distillation-Läufe wurden diese Bestandteile nicht neu konzipiert, sondern als gemeinsame Infrastruktur übernommen, damit der Effekt von Teacher, Student sowie Temperatur und $\alpha$ isoliert untersucht werden kann.

\subsection{Evaluationsmetrik}
\label{subsec:evaluation}

Zur Bewertung der Klassifikationsleistung wird die \textit{Accuracy} (Genauigkeit) herangezogen, definiert als der Anteil korrekt klassifizierter Bilder an der Gesamtzahl der Testbilder:
\begin{equation}
  \text{Accuracy} = \frac{\text{Anzahl korrekt klassifizierter Bilder}}{\text{Gesamtzahl der Testbilder}}
\end{equation}

Da der PatternNet-Datensatz gleichverteilte Klassen aufweist, ist die Accuracy eine aussagekräftige Metrik, die nicht durch Klassenungleichgewichte verzerrt wird.
Darüber hinaus wurde die Konfusionsmatrix verwendet, um die Verteilung der Fehlklassifikationen über die Klassen hinweg zu analysieren und potenzielle Schwächen des Modells in Bezug auf bestimmte Klassen zu identifizieren.